% Options for packages loaded elsewhere
\PassOptionsToPackage{unicode}{hyperref}
\PassOptionsToPackage{hyphens}{url}
\documentclass[
]{article}
\usepackage{xcolor}
\usepackage{amsmath,amssymb}
\setcounter{secnumdepth}{-\maxdimen} % remove section numbering
\usepackage{iftex}
\ifPDFTeX
  \usepackage[T1]{fontenc}
  \usepackage[utf8]{inputenc}
  \usepackage{textcomp} % provide euro and other symbols
\else % if luatex or xetex
  \usepackage{unicode-math} % this also loads fontspec
  \defaultfontfeatures{Scale=MatchLowercase}
  \defaultfontfeatures[\rmfamily]{Ligatures=TeX,Scale=1}
\fi
\usepackage{lmodern}
\ifPDFTeX\else
  % xetex/luatex font selection
\fi
% Use upquote if available, for straight quotes in verbatim environments
\IfFileExists{upquote.sty}{\usepackage{upquote}}{}
\IfFileExists{microtype.sty}{% use microtype if available
  \usepackage[]{microtype}
  \UseMicrotypeSet[protrusion]{basicmath} % disable protrusion for tt fonts
}{}
\makeatletter
\@ifundefined{KOMAClassName}{% if non-KOMA class
  \IfFileExists{parskip.sty}{%
    \usepackage{parskip}
  }{% else
    \setlength{\parindent}{0pt}
    \setlength{\parskip}{6pt plus 2pt minus 1pt}}
}{% if KOMA class
  \KOMAoptions{parskip=half}}
\makeatother
\usepackage{longtable,booktabs,array}
\usepackage{calc} % for calculating minipage widths
% Correct order of tables after \paragraph or \subparagraph
\usepackage{etoolbox}
\makeatletter
\patchcmd\longtable{\par}{\if@noskipsec\mbox{}\fi\par}{}{}
\makeatother
% Allow footnotes in longtable head/foot
\IfFileExists{footnotehyper.sty}{\usepackage{footnotehyper}}{\usepackage{footnote}}
\makesavenoteenv{longtable}
\setlength{\emergencystretch}{3em} % prevent overfull lines
\providecommand{\tightlist}{%
  \setlength{\itemsep}{0pt}\setlength{\parskip}{0pt}}
\usepackage{bookmark}
\IfFileExists{xurl.sty}{\usepackage{xurl}}{} % add URL line breaks if available
\urlstyle{same}
\hypersetup{
  hidelinks,
  pdfcreator={LaTeX via pandoc}}

\author{}
\date{}

\begin{document}

\section{论文草稿 P1 (完整版) --- Generalization as Formalized
Induction}\label{ux8bbaux6587ux8349ux7a3f-p1-ux5b8cux6574ux7248-generalization-as-formalized-induction}

\begin{quote}
版本: v1.0 草稿 \textbar{} 日期: 2026-08-11 结构: 用户大纲 10 节 +
Appendix。理论命题细节见 论文草稿\_P1\_泛化作为归纳形式化.md (§1.0-1.7)
与 结构直觉二象性研究.md。 语言: 标题/关键命题英文, 论证中文 (工作稿;
投稿前全英文)。 纪律: 效果数据双口径 (判定口径 = 主锚)。标注
{[}确认中{]} 的数值需归档复核。
\end{quote}

\begin{center}\rule{0.5\linewidth}{0.5pt}\end{center}

\subsection{Title (方向)}\label{title-ux65b9ux5411}

\emph{Generalization as Formalized Induction: Token-Native Factorized
Representation and Weight Compilation in a Minimal Transformer}

\subsection{Abstract (草稿)}\label{abstract-ux8349ux7a3f}

We propose that systematic generalization is not an emergent or a priori
faculty of neural networks, but the formalization of induction. To
generalize on a concept A, A must be precisely defined, and every
concept in A's definition chain must be sufficiently trained; sufficient
training is a relational property --- parallel concepts under the same
superordinate concept, wrong-answer examples, and collocation coverage
--- not a quantitative one. We operationalize this in a token-native
representation (B/C/S/G/P/A projection layers) in which conceptual
degrees of freedom are factorized according to their symmetries, and
execution is driven by token definitions (a zero-hardcoding evaluator).
A 2-layer, \textless1 MB transformer trained on \textasciitilde2000
samples in a few epochs generalizes to arbitrary width (2000-digit
addition, acc 1.000), unseen radices ({[}3..9{]}, OOD 0.998), unseen
operator orders (iteration chain), and joint Cartesian OOD. Symbol
permutation invariance holds (0.996→0.987 under 26080 digit renamings),
and the three execution channels --- construction, intuition, formal
rewriting --- are semantically equivalent (1.000 on 2000-digit carry
chains), the neural implementation of Computational Trinitarianism. We
emphasize that intuition is the reasoning fast path: it is hard to
obtain (acquisition requires config-dependent epochs --- 3 for simple, 8
for the full logic+arithmetic suite), but once formed it generalizes
easily by skipping steps --- jumping from input to output over the
intermediate construction steps. Acquisition and step-skipping
generalization are decoupled.

\begin{center}\rule{0.5\linewidth}{0.5pt}\end{center}

\subsection{1. Introduction}\label{introduction}

\subsubsection{1.1
现有问题是什么}\label{ux73b0ux6709ux95eeux9898ux662fux4ec0ux4e48}

\begin{enumerate}
\def\labelenumi{\arabic{enumi}.}
\tightlist
\item
  \textbf{系统化泛化是记录的现象, 没有发生学解释}。Compositional
  generalization (Lake et al., 2017; Keysers et al., 2020 COGS; SCAN)
  被广泛观测与基准化, 但文献回答''模型能不能泛化'',
  不回答''泛化如何产生''。系统化泛化被当作模型的涌现属性或数据工程的结果。
\item
  \textbf{主流路线预设 primitive, 跳过概念形成}。λ-calculus
  及''已有数学语言 → tokenize → transformer''路线
  (application/abstraction/β-reduction 作训练目标) 把 primitive
  当作已给: Church 编码可以\emph{实现}一切, 但编码抹平结构边界 ---
  泛化所需的结构 (自由度、对称性、定义链) 在表示中不可见,
  模型只能记忆编码后的函数表。这解释了为何此类模型长度/进制/阶数外推脆弱。
\item
  \textbf{两个对立面都不成立}:

  \begin{itemize}
  \tightlist
  \item
    ``泛化是先验/涌现能力'' (天赋论) --- 无法工程化, 且被我们实验否定:
    定义模糊的 token 即使大量样本也不泛化 (E0/E61);
  \item
    ``泛化只是 benchmark engineering / 表示工程 + 大量数据'' ---
    本工作证明表示工程是\emph{正当贡献}: 当表示因子化后,
    泛化是极低复杂度的学习问题 (结构搜索由表示设计完成, SGD
    只需保持结构)。
  \end{itemize}
\item
  \textbf{神经符号/算术模块的局限} (NALU, Neural Arithmetic Logic
  Units): 固定算子集, 一次训练只能学会预先枚举的算子;
  无法应对''训练后新增算子/新自由度组合''。
\item
  \textbf{结果被低估的风险}: 极小模型 (2层/\textless1MB/2000样本/3epoch)
  全泛化, 若只报数字, 会被质疑泄漏/运气。本论文主动把
  \emph{representation engineering 本身作为贡献}, 使攻击面最小。
\end{enumerate}

\subsubsection{1.2 Babel 做了什么}\label{babel-ux505aux4e86ux4ec0ux4e48}

Babel 是一个 token 原生的表示与执行系统, 基于一个简单主张:
\textbf{泛化是归纳的形式化} --- 对概念 A 泛化 ⟺ ①精确形式化定义 A, ②定义
A 所需的全部 token 充分训练。由此:

\begin{itemize}
\tightlist
\item
  \textbf{表示}: 概念自由度按对称性因子化 (base × direction × inverse ×
  level × domain), 六层投影 (B 公设/C 概念/S 符号/G 语法/P 呈现/A 箭头)
  使能指与所指在 token 层分离。
\item
  \textbf{执行}: 定义驱动的通用求值器 (零硬编码), 真值表/归约规则从
  token 定义读取; 提供构造通道 (展开执行)。
\item
  \textbf{训练}: 训练样本 = 定义的实例化 (平行概念 + 错题 + 搭配),
  监督只含判定真值、不含展开 trace --- 模型必须自己把构造编译进权重。
\item
  \textbf{结果}: 极小模型系统化泛化 (任意位宽/进制/阶数/组合),
  符号置换不变性 (0.987), 三通道语义等价 (1.000), 语义等价路径编译
  (推理快路径)。
\end{itemize}

\subsubsection{1.3 核心贡献 (5 条, 每条 =
可证伪命题)}\label{ux6838ux5fc3ux8d21ux732e-5-ux6761-ux6bcfux6761-ux53efux8bc1ux4f2aux547dux9898}

\begin{itemize}
\tightlist
\item
  \textbf{C1 理论 --- 泛化 = 归纳的形式化 (Formalized Induction)}:
  系统化泛化不是涌现/先验, 是可工程化的: 定义完备性 (G1) + 关系充分性
  (G2) 是泛化的充分条件。反先验论证经验可证伪 (E0, E61)。
\item
  \textbf{C2 系统 --- token 原生因子化表示}: 六层投影 (B/C/S/G/P/A) +
  自由度因子化 + 零硬编码求值引擎, 使''定义精确''可验证
  (无孤儿引用/环证书/覆盖审计), 使表示本身成为可训练、可验证的对象。
\item
  \textbf{C3 实证 --- 极小模型系统化泛化}: 2 层 \textless1MB
  \textasciitilde2000 样本, 任意位宽 (2000 位) / unseen radix / unseen
  order / 联合笛卡尔 OOD; 符号置换不变性 (0.987) + 跳步泛化零衰减 (2000
  位/进制60/100\%匿名/联合 全 1.000)。
\item
  \textbf{C4 机制 --- 权重编译 (Weight Compilation)}: 训练 =
  把定义驱动展开 (构造通道) 编译进权重; 监督不含 answer trace,
  故模型学的是结构而非复制展开过程。产物是\textbf{直觉 --- 推理快路径}:
  难以获得 (需充分训练), 但一旦形成\textbf{容易跳步泛化}
  (跳过构造中间步骤直达结果)。
\item
  \textbf{C5 认识论 --- 三通道等价 (Three-Channel Equivalence)}:
  构造/直觉/形式三通道对同一输入语义等价 (exp50: 2000 位进位链 1.000)
  --- Computational Trinitarianism (Curry-Howard) 的神经实现版,
  等价发生在同一系统执行通道间而非逻辑间。
\end{itemize}

\begin{center}\rule{0.5\linewidth}{0.5pt}\end{center}

\subsection{2. Problem Setting}\label{problem-setting}

\subsubsection{2.1 什么叫 systematic
generalization}\label{ux4ec0ux4e48ux53eb-systematic-generalization}

\begin{itemize}
\tightlist
\item
  定义 (沿用组合泛化文献):
  模型对训练中未出现的\textbf{组合}正确响应的能力。组合发生在独立自由度之间:
  位宽 × 进制 × 算符阶 × 定义域 × 嵌套深度。
\item
  与''插值''区分: 本文测试均为训练分布之外的组合
  (未出现的位宽/进制/阶/笛卡尔组合), 非插值。
\end{itemize}

\subsubsection{2.2 什么是训练内 /
OOD}\label{ux4ec0ux4e48ux662fux8badux7ec3ux5185-ood}

\begin{itemize}
\tightlist
\item
  \textbf{训练内 (in-distribution)}:
  该自由度的取值组合出现在训练样本中。
\item
  \textbf{OOD (out-of-distribution)}: 单自由度外推 (位宽超训练长度;
  进制不在训练集; 算符阶越过训练最高阶) 或联合外推
  (多个自由度同时取未见组合 --- 笛卡尔积中未采样的格点)。
\item
  判定口径: 完整输出序列重建的末尾真值 acc (主锚); 位置级 acc (每位置逐
  token 正确率) 只作参考 --- 位置级存在''digit\_nine 0.16 位置假象 vs
  判定 1.000''类误导, 不可单独作为结论。
\end{itemize}

\subsubsection{2.3 Babel 支持的
grammar}\label{babel-ux652fux6301ux7684-grammar}

\begin{itemize}
\tightlist
\item
  表达式文法由 token 数据驱动 (G 层 gtoken 排列 + P 层呈现), 槽位角色:
  arg/fn/args/binder/body。节点类型: atom / application / equality /
  binary\_connective / unary\_connective / quantified / eos。
\item
  numeral: {[}sign + 进制部分{]} + digit 序列 (内禀特征向量 × digit
  向量), 任意进制 (无单字符名 digit 亦可)。
\item
  prop: 逻辑命题判定序列 (judgment), 判定格式沿 gtoken/ptoken 组装,
  零手写。
\item
  覆盖: 算术 (加减乘除幂根/超运算/neg), 逻辑 (9 门), 比较, 迭代链,
  数域提升 (coercion)。
\end{itemize}

\begin{center}\rule{0.5\linewidth}{0.5pt}\end{center}

\subsection{3. Babel (表示系统)}\label{babel-ux8868ux793aux7cfbux7edf}

\subsubsection{3.1 六层投影 (token 是''一'', 多 json
是投影)}\label{ux516dux5c42ux6295ux5f71-token-ux662fux4e00-ux591a-json-ux662fux6295ux5f71}

\begin{longtable}[]{@{}
  >{\raggedright\arraybackslash}p{(\linewidth - 4\tabcolsep) * \real{0.3333}}
  >{\raggedright\arraybackslash}p{(\linewidth - 4\tabcolsep) * \real{0.3333}}
  >{\raggedright\arraybackslash}p{(\linewidth - 4\tabcolsep) * \real{0.3333}}@{}}
\toprule\noalign{}
\begin{minipage}[b]{\linewidth}\raggedright
层
\end{minipage} & \begin{minipage}[b]{\linewidth}\raggedright
实体
\end{minipage} & \begin{minipage}[b]{\linewidth}\raggedright
作用
\end{minipage} \\
\midrule\noalign{}
\endhead
\bottomrule\noalign{}
\endlastfoot
B 公设基 & baseloop.jsonl & axiomatic 终点, 定义链溯源终点 (brace
穿透终止) \\
C 概念 & derive.jsonl & 5 类定义形态
(explicit/inductive/recursive/axiomatic/implicit) \\
S 符号 & symbol\_tokens.jsonl & 能指 (glyph → maps\_to 多对多,
歧义保留全候选, 上下文收敛) \\
G 语法 & gtokens.jsonl & 排列方法 (槽位角色/归约), 语法本身是 token \\
P 呈现 & presentation.jsonl & 具体记法 (符号/优先级/结合性),
一概念多呈现 \\
A 箭头 & arrow\_tokens.jsonl & 概念间有向关系
(数域提升/迭代链/对偶/coercion) \\
\end{longtable}

核心设计: 变换只依赖 token 本身; 知识在 token/语法数据里, 代码只读取执行
(零硬编码纪律)。

\subsubsection{3.2 自由度}\label{ux81eaux7531ux5ea6}

概念 = 独立自由度坐标, 而非孤立 token: A = (d1, \ldots, dk)。例: -
迭代链: 阶 (succ→+→×→\^{}→↑↑) --- 自由度''层'' - 对偶: inverse
(succ/pred, power/root, +/−) --- 自由度''方向'' - 数域:
natural/integer/rational/real/complex --- 自由度''domain'' - 表示法:
sign/base/cardinality --- 自由度''表示''

\subsubsection{3.3 对称性}\label{ux5bf9ux79f0ux6027}

\begin{itemize}
\tightlist
\item
  迭代对称: iterate(方向, 层数, 被迭代算符) --- 高阶算符;
  F\_\{n+1\}(a,b) = Iterate(F\_n, a, b)。
\item
  对偶对称: inverse 箭头 (power↔root 等), 对称家族 (neg/scale/root;
  complement/parallel\_sum De Morgan 对偶; translation/inversion 模群)。
\item
  数域提升: coercion 单态射 (A 层 concept=coercion 检索, 零硬编码)。
\end{itemize}

\subsubsection{3.4 composition}\label{composition}

\begin{itemize}
\tightlist
\item
  排列方法在 G 层 (gtoken 槽位), 组装器只读数据; SKI 组合子亦为 gtoken
  (归约规则在 reduction 数据)。
\item
  组装: 概念 → arrange → gtoken → 嵌套 AST → 呈现文本 (数据驱动,
  parse/print 双向)。
\end{itemize}

\subsubsection{3.5 iteration}\label{iteration}

\begin{itemize}
\tightlist
\item
  iterate 是被迭代算符上的算符: {[}方向, fn, 层数,
  被迭代算符{]}。迭代链成员沿 A 层箭头发现 (ARROW\_REGISTRY), 零
  is\_succ 硬编码混用。
\item
  通用求值引擎: 纯迭代链成员走迭代归约, 非链成员
  (translation/inversion/differential 等) 委托对称求值 --- token 进
  token 出, eid 索引, 零每次查名。
\end{itemize}

\subsubsection{3.6 domain/radix
representation}\label{domainradix-representation}

\begin{itemize}
\tightlist
\item
  进制是\emph{对象}不是结构常数: numeral = {[}sign, base, cardinality{]}
  ⊗ digit 向量; 同一组权重处理任意 base (base 作为输入参数,
  不固化进算子)。
\item
  表示法可扩展 (fraction/radical/imaginary\_unit), 数符序列保持 numeral
  (digit 排列), 禁''直接按 value''丢 digit。
\item
  包裹分离: 操作嵌套 ( ) vs 表示嵌套 {[} {]} 必须不同符号 --- 同符号 =
  模型分不清运算分组 vs 数表示结构 (训练崩, 模态分离 E4.1)。
\end{itemize}

\begin{center}\rule{0.5\linewidth}{0.5pt}\end{center}

\subsection{4. Weight Compilation}\label{weight-compilation}

\subsubsection{4.1
数据如何生成}\label{ux6570ux636eux5982ux4f55ux751fux6210}

\begin{itemize}
\tightlist
\item
  定义驱动: 样本 = 定义的实例化。合成器 (synth\_core, 配置驱动) 组合:
  balanced/trace/numeral\_split/fill/choose/definition/arrow/law/logic\_nested/logic\_arith/logic\_interdef/coercion/deep\_nest/extrap/cartesian/radix/neg。
\item
  充分训练三要素 (G2): 平行概念同训 (关系成为学习信号) + 错题 (负例,
  含同构结果错) + 搭配 (相邻搭配对 ≥3, 低频搭配重分配)。
\item
  零手写序列: 拼装沿 gtoken/ptoken (assemble\_seq/judge\_seq), 手写序列
  = ``缺 token''信号。
\end{itemize}

\subsubsection{4.2 Python reference
evaluator}\label{python-reference-evaluator}

\begin{itemize}
\tightlist
\item
  通用求值引擎 (eval/engine.py): eval\_op(op, arg\_tokens) →
  result\_tokens, token 进 token 出。
\item
  真值由 token 定义提供: 逻辑真值表/比较真值表/归约规则全部读定义
  (logic\_eval/compare\_eval/digit\_eval/arith\_eval/symmetry\_eval/prime\_eval/token\_iterator)。
\item
  双语义来源一致: 构造通道 (token\_iterator 展开) 与形式通道 (引擎重写)
  同源读定义 → 三通道一致性的基础。
\end{itemize}

\subsubsection{4.3 Transformer
如何训练}\label{transformer-ux5982ux4f55ux8badux7ec3}

\begin{itemize}
\tightlist
\item
  极小架构: 2 层 transformer, dim 64, \textless1MB 权重。
\item
  目标: 位置 CE (正例) + 合法性 CE (对错区分)。
\item
  训练量: \textasciitilde2000 样本, 收敛 epoch 配置相关 (简单配置 3
  epochs, 完整逻辑+算术套件 8 epochs --- exp53 曲线)。直觉难以获得,
  但容易跳步泛化 (EXP-80 验证)。
\item
  判定口径 = 完整序列重建 (末尾真值), 位置级只作参考。
\item
  诊断: per\_token\_gen 逐 token 泛化 / epoch\_gen 曲线 / crash\_culprit
  带崩元凶 / 覆盖审计。
\end{itemize}

\subsubsection{4.4 为什么 input 不包含 answer
trace}\label{ux4e3aux4ec0ux4e48-input-ux4e0dux5305ux542b-answer-trace}

\begin{itemize}
\tightlist
\item
  监督只含判定真值 (正/误), \textbf{不含逐步展开的
  trace}。这是''编译''而非''蒸馏''的关键: 若给 trace,
  模型只需复制展开过程 (序列克隆); 不给 trace,
  模型必须\textbf{自己从定义归纳出计算结构并编译进权重} --- 学习即编译。
\item
  对应三通道: 训练 = 构造 (展开执行, oracle 生成监督) → 编译 → 直觉
  (推理快路径, 跳步直达)。
\item
  这使快路径具备真正的压缩: 不是记忆展开, 而是结构编译 (weights ≈
  program, tokens ≈ tape)。
\end{itemize}

\begin{center}\rule{0.5\linewidth}{0.5pt}\end{center}

\subsection{5. Experiments}\label{experiments}

\begin{quote}
全部走 run\_exp 统一归档 (五件套:
config/metrics/model/samples/views)。双口径: 判定 acc / 位置级
acc。{[}确认中{]} = 归档复核数值。
\end{quote}

\begin{longtable}[]{@{}
  >{\raggedright\arraybackslash}p{(\linewidth - 4\tabcolsep) * \real{0.3333}}
  >{\raggedright\arraybackslash}p{(\linewidth - 4\tabcolsep) * \real{0.3333}}
  >{\raggedright\arraybackslash}p{(\linewidth - 4\tabcolsep) * \real{0.3333}}@{}}
\toprule\noalign{}
\begin{minipage}[b]{\linewidth}\raggedright
实验
\end{minipage} & \begin{minipage}[b]{\linewidth}\raggedright
结果 (判定口径)
\end{minipage} & \begin{minipage}[b]{\linewidth}\raggedright
状态
\end{minipage} \\
\midrule\noalign{}
\endhead
\bottomrule\noalign{}
\endlastfoot
基线: known operators (8+ 算术/逻辑门) & arith\_multi\_type 判定 1.000
(1188 样本) & ✅ \\
2000-digit addition 外推 & acc 1.000 & ✅ \\
unseen radix (进制 {[}3..9{]} × 20 位) & OOD 0.998 & ✅ \\
unseen order (迭代链越阶) & 全泛化 & ✅ \\
joint Cartesian OOD (位宽×进制×阶×嵌套) & 判定 {[}确认中{]} (0.816 链 /
0.988 组) & ✅/复核 \\
\textbf{符号置换 (exp41)} & 26080 处 digit 重命名, 0.996→0.987
(置换不变性成立) & ✅ \\
\textbf{三通道一致性 (exp50)} & 构造=形式 (2/9/20/2000 位), 直觉 1.000;
2000 位进位链三通道全对 & ✅ \\
\textbf{快慢成本 (exp51)} & 2000 位: 构造 0.62s / 直觉 0.086s (7×) /
形式 0.069s (9×) & ✅ \\
\textbf{平行概念移出 (exp20)} & 删 iff/xor → 0.798 (其他门 1.0→0.833);
删 and/or → 0.814 & ✅ \\
\textbf{收敛曲线 (exp53)} & epoch 1/2/3/5/8 →
0.001/0.214/0.244/0.781/0.996 --- SUPPORT: 直觉难以获得, 收敛点配置相关
(简单 3 / 完整 8) & ✅ \\
\textbf{跳步泛化零衰减 (exp80)} & 2000 位 / 进制 60 / 100\% digit 匿名 /
位宽×进制联合 \textbf{全 1.000} (获得 8 epochs 解耦) & ✅ \\
\textbf{跳步机器形态 (exp90, 校准版)} & 20 位加法: 结构 82 token/0.410ms
vs 直觉 5 token/0.186ms (num 原子化, 16× token 缩短, 2.2×), 语义一致;
口径 = 序列长度效应下限 (itoken 不在 vocab) & ✅ \\
\textbf{样本量 (exp60)} & 随机裁剪破搭配覆盖, 测法无效 --- 不作为 G6
证据 (待确定性裁剪重测) & 方法层 \\
\textbf{imply 判定监督 (exp10\_imply\_supervised)} &
定义+位置+判定监督三要素 ⇒ 判定 0.996 (imply 与 9 门全部学会) & ✅ \\
anonymous operators (符号置换全 operator 版) & {[}待做{]} & ⏳ \\
模型大小/样本/epoch & 2 层 / 64 dim / \textless1MB / \textasciitilde2000
/ 3 epochs & ✅ \\
\end{longtable}

\subsubsection{5.x
关键实验要点}\label{x-ux5173ux952eux5b9eux9a8cux8981ux70b9}

\begin{itemize}
\tightlist
\item
  \textbf{known operators}: 加减乘除幂根/超运算/neg + 逻辑门,
  统一参数集, 非逐算子训练。
\item
  \textbf{2000-digit}: 训练位宽 ≤ 有限 L, 测试 2000 位 --- 长度无关,
  砍''函数表记忆''。
\item
  \textbf{unseen radix}: base 是输入参数, 非结构常数 ---
  砍''十进制记忆''。
\item
  \textbf{unseen domain/order}: 越过训练最高阶泛化 --- 学到算符生成规则
  (迭代), 非逐算子。
\item
  \textbf{joint OOD}: 未出现自由度组合, 判定口径全序列重建。
\item
  \textbf{imply 判定监督 (exp10\_imply\_supervised)}:
  定义+位置+判定监督三要素齐全时 imply 与 9 门全部学会 (判定
  0.996)。imply 需自身判定监督 (exp20 两变体 imply 均 0 即证) ---
  与''定义等式传递语义, 判定形式需位置+监督覆盖''一致。
\item
  \textbf{anonymous operators}: digit/operator 全随机重命名, 保持结构
  --- 学的是关系非符号语义 (exp41 置换 0.987 / exp80 100\% 匿名 1.000)。
\end{itemize}

\begin{center}\rule{0.5\linewidth}{0.5pt}\end{center}

\subsection{6. 语义等价路径编译 (Semantic-Path
Compilation)}\label{ux8bedux4e49ux7b49ux4ef7ux8defux5f84ux7f16ux8bd1-semantic-path-compilation}

\subsubsection{6.1 0
样本已经能执行}\label{ux6837ux672cux5df2ux7ecfux80fdux6267ux884c}

\begin{itemize}
\tightlist
\item
  新算符 σ 的语义闭包由定义链保证: 基础模型 K0 未见 σ, 但沿展开路径 T\_σ
  仍可执行 K0(T\_σ, x) = y --- 只是路径长 (构造通道)。
\item
  即: \textbf{语义闭包不依赖训练}; 0 样本时已有正确慢执行路径 (exp50:
  构造通道 2000 位进位链正确)。
\end{itemize}

\subsubsection{6.2 少样本建立推理快路径 --- 直觉难以获得, 容易跳步泛化
(术语定稿)}\label{ux5c11ux6837ux672cux5efaux7acbux63a8ux7406ux5febux8defux5f84-ux76f4ux89c9ux96beux4ee5ux83b7ux5f97-ux5bb9ux6613ux8df3ux6b65ux6cdbux5316-ux672fux8bedux5b9aux7a3f}

\begin{itemize}
\tightlist
\item
  \textbf{术语 (发明人定稿)}: 直觉 =
  \textbf{推理快路径}。它\textbf{难以获得} (训练需定义完备 + 关系充分,
  收敛 epoch 配置相关 --- exp53 完整套件 8 epochs),
  但一旦形成\textbf{容易跳步泛化} ---
  泛化时跳过构造中间步骤直达结果。快是''推理快路径''的本质,
  不是独立优化; 跳步是它对外泛化容易的机制。
\item
  少样本建立语义等价直觉通道: 从构造通道自动生成监督, 编译出
  Sem\_fast(σ) = Sem\_expanded(T\_σ) (可验证/fallback); Cost\_fast
  \textless{} Cost\_expanded (exp51: 2000 位 7×)。
\item
  直觉通道不扩大语义闭包 --- 解决执行成本问题, 不是学习新语义。
\item
  与经典蒸馏区别: teacher 与 student
  是\textbf{同一模型}的构造/直觉两条通道 (构造通道自蒸馏);
  语义等价有验证器与 fallback。
\item
  诚实声明: 形式引擎 (eval\_op, 0.069s) 比直觉 (0.086s) 快 ---
  直觉的价值是\textbf{端到端} (权重内执行, 无需外部 oracle) +
  \textbf{跳步泛化保持} (exp41 置换 0.987 / exp50 2000 位 1.000),
  非绝对速度。
\end{itemize}

\begin{center}\rule{0.5\linewidth}{0.5pt}\end{center}

\subsection{7. Ablations}\label{ablations}

\begin{longtable}[]{@{}
  >{\raggedright\arraybackslash}p{(\linewidth - 4\tabcolsep) * \real{0.3333}}
  >{\raggedright\arraybackslash}p{(\linewidth - 4\tabcolsep) * \real{0.3333}}
  >{\raggedright\arraybackslash}p{(\linewidth - 4\tabcolsep) * \real{0.3333}}@{}}
\toprule\noalign{}
\begin{minipage}[b]{\linewidth}\raggedright
Ablation
\end{minipage} & \begin{minipage}[b]{\linewidth}\raggedright
结果
\end{minipage} & \begin{minipage}[b]{\linewidth}\raggedright
支持命题
\end{minipage} \\
\midrule\noalign{}
\endhead
\bottomrule\noalign{}
\endlastfoot
拆自由度 (numeral\_v4 正确) vs 纠缠/错对称 (wrong-symmetry/ablated) &
纠缠 ⇒ 泛化崩 & G4/C1 \\
\textbf{平行概念移出 (exp20)} & 删 iff/xor → 整体 0.798, 被删门 0.000,
其他门 1.0→0.833; 删 and/or → 0.814 & G2 \\
\textbf{符号置换 (exp41)} & 26080 处 digit 重命名 0.996→0.987
(置换不变性) & G4 \\
\textbf{三通道一致性 (exp50)} & 构造=形式=直觉 (2000 位进位链 1.000) &
R/C5 \\
\textbf{收敛曲线 (exp53)} & 完整配置 1/2/3/5/8 →
0.001/0.214/0.244/0.781/0.996 --- SUPPORT: 直觉难以获得 & G5 \\
\textbf{跳步泛化解耦 (exp80)} & 2000 位/进制 60/100\% 匿名/联合全 1.000,
获得 8 epochs --- 难以获得, 容易跳步泛化 & G5/R \\
\textbf{样本量 (exp60)} & 随机裁剪破搭配, 测法无效 --- 不作证据
(待确定性裁剪) & 方法层 \\
负例移除 (neg 覆盖 9→15) & neg 0.000→0.958 & G2 \\
定义链 token 未覆盖 (value\_three 样本量=1) & 学不会, 覆盖后学会 & G1 \\
定义模糊 (vocab 153) & 需 40 epochs 才收敛 & G0/G6 \\
包裹符号混用 ( ( ) vs {[} {]} 同符号) & 训练崩 & G4 \\
判定口径 vs 位置级 (focused) & 位置 0.904 → 判定 0.000 (假象) &
口径纪律 \\
\textbf{普通 representation baseline} (EXP-C1: radix-fixed /
iteration-hidden / entangled) & {[}待做{]} & C1 \\
\end{longtable}

普通 baseline (EXP-C1) 是决定性对照: 相同模型/数据预算, 表示去因子化
(进制固化/迭代隐藏/纠缠编码) ⇒ 外推崩 (≤0.5); 因子化参照 (exp02)
外推零衰减 (1.000)。证明因子化表示是泛化的\emph{原因}非巧合。

\begin{quote}
\textbf{方法纪律 (§7.2 教训)}: EXP-10 (0-IMPLY) 系列因评估口径 bug
全部撤回 --- 原''零监督语义涌现''为末尾真值口径 + 语法错位假象。教训:
判定口径 (全序列重建) 唯一可信; 评估 branch 必须走 run\_exp 官方路径;
位置覆盖是独立自由度 (imply 需判定监督 + 位置覆盖,
exp10\_imply\_supervised 0.996 即证三要素)。不构成本文证据,
构成方法纪律。
\end{quote}

\begin{center}\rule{0.5\linewidth}{0.5pt}\end{center}

\subsection{8. Related Work}\label{related-work}

\begin{enumerate}
\def\labelenumi{\arabic{enumi}.}
\tightlist
\item
  \textbf{λ-calculus 与形式系统家族} --- 作为被解释对象: 预设 primitive,
  不回答概念形成; Church 编码抹平结构。本文提供''表示的形式化'':
  能指/所指分离、歧义为一等公民、排列方法可定义、抽象/具体记法分离、公设基分层、概念网络、自由度分解。
\item
  \textbf{系统化泛化} --- Lake et al.~2017; Keysers et al.~2020 (COGS);
  SCAN; 组合性综述: 记录现象; 本文提供发生学解释 (为何/何时泛化成立)。
\item
  \textbf{神经符号与算术模块} --- NALU/NAU (固定算子集), Neural
  Arithmetic: 一次一算子; 本文统一参数集 + 新增算子零样本语义闭包。
\item
  \textbf{长度泛化} --- 位置编码/循环消解系列: 固定表示下的工程对策;
  本文从表示层消除长度相关性。
\item
  \textbf{表示学习} --- disentangled/compositional structured
  representations: 研究表示''长什么样'';
  本文回答''表示如何决定可学习性'', 且把表示设计作为贡献而非工具。
\item
  \textbf{程序 → 神经编译} --- Learning to Compile Programs to Neural
  Networks (ICML 2024) 等: 相似方法面; 区别 = 同一模型慢路径自蒸馏,
  不扩语义闭包, 可验证+fallback。
\item
  \textbf{符号学与认知双系统} --- Frege (Sinn/Bedeutung), Saussure
  (能指/所指/任意性/差异原则): G4 祖先; Kahneman (System 1/2): G5 对应。
\item
  \textbf{数学哲学 --- 三通道等价站位} --- 数学结构本身全部已有 (de
  Bruijn/OpenMath/组合子逻辑/Bourbaki/范畴论/Mizar-Lean-Isabelle),
  \textbf{不 claim}; 只 claim 三通道等价
  (构造/直觉/形式在同一系统执行通道间可编译等价) = Computational
  Trinitarianism 神经实现版。立场: 反柏拉图主义 (直觉 = 编译的构造,
  非天赋)。
\item
  \textbf{反天赋论} --- Fodor (language of thought 天赋概念), Chomsky
  (普遍语法): 对立面; 经验反证 E0/E61。
\end{enumerate}

\begin{center}\rule{0.5\linewidth}{0.5pt}\end{center}

\subsection{9. Limitations}\label{limitations}

\begin{enumerate}
\def\labelenumi{\arabic{enumi}.}
\tightlist
\item
  \textbf{serialization (位置模式偏)}: 位置覆盖是独立于语义的自由度 ---
  序列中某语法槽位 (如判定首 token) 未被训练覆盖时,
  即使定义完备也不泛化; 需位置覆盖训练。位置依赖在 AR
  生成中不可完全消除。
\item
  \textbf{finite resources}: 权重/样本/epoch 存在下界; 收敛 epoch
  配置相关 (简单 3, 完整 8 --- exp53); 超运算结果天文数字,
  判定型负例用小值域 + max\_val 防溢出 (不声称无限精度)。
\item
  \textbf{supported calculus ≠ 无限数学}: 系统只覆盖已形式化定义的结构
  (定义形态内); 不声称解决一般数学或自主概念发现;
  新概念仍需人工登记定义。
\item
  \textbf{值空间/表示}: 当前主要做 bool 判定与计数结构, 完整值空间
  (分数/复数) 部分接入, numeral\_value 侧未全部开放。
\item
  \textbf{直觉不是保证性的}: 直觉通道是统计性结构定位, 非逻辑保证 ---
  需验证器/fallback (构造通道) 背书等价性。
\end{enumerate}

\begin{center}\rule{0.5\linewidth}{0.5pt}\end{center}

\subsection{10. Conclusion}\label{conclusion}

\begin{itemize}
\tightlist
\item
  泛化 = 归纳的形式化: 定义完备 + 关系充分 ⇒ 极小模型系统化泛化,
  可工程化, 可复现。
\item
  表示是贡献: 自由度因子化使泛化成为低复杂度学习问题 ---
  把泛化写进表示的几何结构, 而非期待模型自行发现。
\item
  权重编译: 训练 = 把构造编译进直觉; 监督不含 trace,
  学的是结构。\textbf{直觉 = 推理快路径}: 难以获得, 但容易跳步泛化 ---
  获得与跳步泛化解耦。
\item
  三通道等价: 构造/直觉/形式在神经实现中互可编译 (exp50: 2000
  位进位链三通道一致) --- 数学基础百年分歧的工程化统一。
\item
  未来: 更大定义域, 自主概念形成 (从数据归纳新定义), 跨模态
  (符号置换的延伸)。
\end{itemize}

\begin{center}\rule{0.5\linewidth}{0.5pt}\end{center}

\subsection{Appendix}\label{appendix}

\begin{itemize}
\tightlist
\item
  A. 完整实验表: 每实验 config/seed/双口径/per\_token 短板/conclusion
  (docs/paper\_data/results.csv)。
\item
  B. Grammar: gtoken/ptoken 完整列表 (数据驱动来源); 呈现定义示例
  (add\_infix/and\_infix/not\_prefix/succ\_prefix)。
\item
  C. 更多 examples: 训练样本序列样例
  (numeral\_split/law/三通道输入/判定序列)。
\item
  D. Training config: 完整 JSON (exp02 套件: synth.samples 组合 + judge
  + train)。
\item
  E. Lean proofs: 定义链完备性 (无孤儿引用/环证书判定) +
  三通道等价的编译正确性形式化 --- 纪律验证, 非新数学 claim。
\end{itemize}

\begin{center}\rule{0.5\linewidth}{0.5pt}\end{center}

\subsection{附: 待办 (草稿 →
论文)}\label{ux9644-ux5f85ux529e-ux8349ux7a3f-ux8bbaux6587}

\begin{itemize}
\tightlist
\item[$\square$]
  EXP-41 符号置换 (anonymous operators) --- 论文第 5 节必填项
\item[$\square$]
  EXP-50 三通道一致性 --- C5 直接证据
\item[$\square$]
  EXP-51 快慢成本对比 --- 第 6 节量化
\item[$\square$]
  普通 representation baseline (第 7 节关键对照)
\item[$\square$]
  数值归档复核 ({[}确认中{]} 项) + paper\_meta 全量登记
\item[$\square$]
  文献精确引用填充 (Related Work)
\item[$\square$]
  Lean 附录形式化
\end{itemize}

\begin{center}\rule{0.5\linewidth}{0.5pt}\end{center}

\subsection{结论 (统一框架, 2026-08-11
定稿)}\label{ux7ed3ux8bba-ux7edfux4e00ux6846ux67b6-2026-08-11-ux5b9aux7a3f}

当我们在统一的形式化注意力语法结构下, 最终不可避免地、必然地通过神经网络
清晰、直观、可解释地观察到直觉与构造的边界时, 我们就可以有充足的信心:

\begin{itemize}
\tightlist
\item
  \textbf{用统一的注意力形式语言指导逻辑构造} (形式→构造: 注意力定位
  token 间关系; 形式 ↔ 注意力 attention);
\item
  \textbf{用精确的合成 CoT 逻辑构造获得 flashing 直觉} (构造→直觉:
  合成思维链 编译为快速推理; 构造 ↔ 合成思维链 synthetic CoT);
\item
  \textbf{用蒸馏的直觉上下文穷举注意力形式} (直觉→形式:
  蒸馏快路径搜索组合, 再表达于统一形式; 直觉 ↔ 蒸馏的闪念 distilled
  intuition / flashing).
\end{itemize}

形式、构造、直觉在统一的框架下相互适配,
组成再不分彼此的教学、训练、推理的 反馈迭代管线 --- 这也许是我们通往 ASI
的众多路径中, 相对较快的一条捷径。

\end{document}
